\section{Discussion}
\label{sec:discussion}

\subsection{What the two-level model contributes}

The value of the two-level formulation is that it forces the policy question to be asked in the
terms the enterprise will answer it. A composition target is not an instruction to the landscape;
it is an aspiration that the enterprise meets, or misses, by choosing which stands to cut and what
to plant, and the cost of that choice is borne in the schedule. Making the two levels explicit is what allows
the analysis to say whether a policy is achievable at the cut it sets, and at what economic cost,
rather than assuming that area shares of species translate into outcomes.

On the TSA~24 estate this distinction is concrete. The highest even flow the compiled estate
sustains is about \SI{51.9}{\mega\cubic\meter} per period, or
\SI{5.19}{\mega\cubic\meter\per\year}, and the schedule is worth \SI{8779}{million\CAD} of net
present value at \SI{3}{\percent} under current prices, costs, and stumpage. That is the feasible
baseline. Any composition policy is a movement away from it, and its cost can now be measured
rather than asserted.

\subsection{Why the model sustains less than the official cut}

The sustained level is about \SI{25}{\percent} below the official allowable annual cut. We treat
this as a result, not a failure to be tuned away, and we can locate it.

It is not the land base. The model's managed area, \SI{3.06}{\mega\hectare}, matches the
base-case timber harvesting land base in the 2017 review, \SI{3.070}{\mega\hectare}, to within
\SI{0.3}{\percent}. It is not the growth potential either: the area-weighted maximum mean annual
increment of the model's own yield curves is \SI{2.31}{\cubic\meter\per\hectare\per\year},
which on the managed area gives \SI{7.07}{\mega\cubic\meter\per\year}, within
\SI{2.5}{\percent} of the published cut per hectare.

The difference is in the harvest-flow formulation. Our policy holds every period to within
\SI{1}{\percent} of the first, while a timber supply review reports a declining transition
profile: the 2017 determination was \SI{8.35}{\mega\cubic\meter\per\year} for five years,
then \SI{7.35}{\mega\cubic\meter\per\year}, and the 2025 adjustment set it at
\SI{6.92}{\mega\cubic\meter\per\year}. A constant level
cannot use the transition harvest that a declining profile allows, and it must survive the trough
in a managed area whose median age is 104 years and whose age structure is far from regulated over a
200-year horizon. The gap is therefore the cost of the flow rule we imposed, not a defect in the
estate.

The practical consequence is that the model and the policy are not answering exactly the same
question. The model asks what this compiled estate can sustain at a constant level, because that is
the object a composition policy has to be compared on. The cut is a decision for a jurisdiction
that includes areas, species, and inventory the model does not carry, and it is itself a transition
level. Reporting the difference as a stated bridge is more useful than forcing the model to
reproduce a number it cannot support.

\subsection{The discount rate dominates}

Across the documented ranges of the economic constants, one lever governs the result: the discount
rate. Moving from \SI{3}{\percent} to \SI{5}{\percent} removes \SI{4010}{million\CAD} of net
present value, close to half the base case, while the deciduous log price across its full range
moves the result by \SI{26}{million\CAD} and haul distance by \SI{184}{million\CAD}.

This ordering has a policy reading that is easy to miss. The constants that are uncertain in the
data are not the ones that decide the answer. The discount rate is not an empirical constant at
all: it is a statement about how much weight the decision maker gives to future harvests, and it is
therefore a policy choice in the same sense as the allowable cut. A policy analysis that reports a
single net present value at a single rate hides the lever that matters most. The sensitivity table
should be read as part of the result, not as an appendix to it.

\subsection{Composition: what the price sensitivity does and does not show}

The insensitivity of mean value to the deciduous price is a real finding, but it is narrower than
it looks. Deciduous species are \SI{5.4}{\percent} of the managed land base and \SI{5.4}{\percent} of
harvested volume, and their unit margin is \SI{-2.25}{CAD\per\cubic\meter}: at the recorded interior
hardwood price, deciduous timber does not cover its delivered cost. Because the share is small and
the margin is close to zero, moving the price across its plausible range changes the total by
little.

What this does not show is that the composition target is economically immaterial. The composition
surface on TSA~24 has not been computed, and varying the price of a fixed deciduous share is not
the same experiment as varying the share itself. A policy that raises the deciduous share changes
which stands are cut and what is planted, and those changes act through the schedule, not only
through the price of the marginal cubic metre.

At the level the estate sustains, the deciduous share can rise from \SI{5.4}{\percent} to about
\SI{6.2}{\percent} at no cost in total harvest, because the physical supply of aspen runs out at
that point. Beyond it, every point of share is bought with level: an eight percent share implies a
total near \SI{40}{\mega\cubic\meter} per period, a ten percent share near \SI{32}{\mega\cubic\meter}, and
a fifteen percent share near \SI{21}{\mega\cubic\meter}. The economic price channel is almost silent; the
physical volume channel is not.

That asymmetry is the central economic finding. A policy analyst who looks only at prices will
conclude that composition is nearly free; a policy analyst who looks at the schedule will find that
it costs a third of the harvest.

\subsection{Fire as a scenario, and the discipline that follows}

Fire enters this analysis as a spatially explicit process calibrated to the historical perimeter
record of the area, not as a probability of loss applied to a stand. That choice is what makes fuel
structure, composition, and connectivity representable, and it is the main methodological
difference from the portfolio approach of \citet{fuchs_severe_2024}.

It also imposes a discipline on the claims. The calibration's central distribution is close to the
historical record, but its upper tail is thin and weakly identified: the historical 90th percentile
of fire size has a bootstrap interval that spans roughly a factor of three, and the share of fires
above \SI{5000}{\hectare} rests on 13 of 432 events. The simulated tail is also sensitive to the
antecedent fire-weather initialisation, which determines how dry the fuel is when a fire starts.
An absolute statement about future fire size would therefore be a statement about a weather
reconstruction and a sample of thirteen events, not about the forest.

We therefore compare policies against each other on identical fire draws, and report the historical
interval beside every distributional comparison. This is the same logic that
\citet{fuchs_severe_2024} applied to extreme weather: the economically relevant question is how a
policy performs relative to the alternatives under a severe regime, not whether the model
reproduces a particular historical event.

The comparison supports the discipline rather than undermining it. The median, the 90th
percentile, the share above \SI{500}{\hectare}, and the maximum are reproduced inside the
historical interval. The 95th percentile and the extreme-tail shares are low. That is the part of
the distribution the policy is meant to manage, and it is the part the calibration cannot pin down.
Reporting the interval beside every statistic is therefore not a hedge; it is the only honest way to
read a model whose tail is thinner than the record.

\subsection{Relation to the European evidence}

The European result was that diversification raises expected returns and reduces isolated risk,
but that severe and frequent extreme events undermine the economic gain and push the optimal
allocation toward homogenisation and low-investment species. Our framework is built to test
whether that conclusion survives when disturbance is a spatial process on a real estate rather than
a correlated draw on stylized units.

\citet{fuchs_severe_2024} found that diversification raises expected returns and reduces isolated
risk, but that severe and frequent extreme events undermine the economic gain. Our result qualifies
that finding rather than contradicting it. The composition lever here is not a free portfolio
reweighting: the minority species is physically limited to about six percent of harvest, so the
choice is constrained before economics enters. And the tail that would justify paying for
diversification is exactly the part of the distribution the model understates. In the European
setting, the constraint was the correlation of extreme events; in the BC setting, it is the supply
of the minority species and the thinness of the modelled tail. Both point the same way: the case
for diversification is conditional, and the condition is empirical.

\subsection{Limitations}

Seven limitations bound what the paper can claim. First, the analysis covers one timber supply area.
TSA~24 is large and fire-prone, but it is one landscape with one inventory, one set of prices, and
one administrative history; the results are not a provincial assessment. Second, growth is
represented by compiled yield curves rather than by individual-tree simulation, so the response of
growth to composition and climate is not modelled. Third, the policy level sets a landscape-wide
composition share rather than a spatial placement rule, so the question of where to put broadleaf
fuel breaks, and how far that can go before production suffers, is not answered here; it is the
natural next step and requires the policy to act on location, not only on share. Fourth, prices,
costs, and stumpage are point values or documented ranges from a single period, and the deciduous
price in particular is an interior hardwood average that may not represent a future market for
small-diameter material. Fifth, the fire model, while calibrated, has not been validated on an
independent period or region, and its tail is weakly identified; the comparative framing protects
the policy conclusions but not a claim about absolute risk.

\subsection{Implications for policy}

Three implications follow, each conditional on the results above. The allowable cut and the
composition target should be evaluated together, because the enterprise's response to one depends
on the other. Risk attitude belongs in the analysis explicitly, because the composition that looks
best at the mean need not be the one that looks best in the tail. And the discount rate deserves
as much scrutiny as the silvicultural prescription, because it moves the result more than any
biological or market parameter in the model.
